\section{Extension to AI planning systems}

\subsection{The step-log as a planning trace}

The surface planner step-log that achieves $\Runcited = 0\%$ in Scenario~4 has the
structure of a planning trace. The agent received a query, resolved a configuration,
dispatched skills, and produced a plan. The step-log records every intermediate state:
which skills ran, what their outputs were, and what the final plan consists of. When
this trace is injected as RAG context, the language model answers from the trace's
artifact references rather than from world knowledge. $\Runcited = 0\%$ because the
trace transforms the query from open-domain science to closed artifact-reference.

The paper ``Reasoning with Language Model is Planning with World Model''~\cite{arxiv230514992}
explicitly frames LLM reasoning steps as world-state traces and proposes evaluating whether
those states are accurate. The evaluation problem $\Runcited$ addresses is the evidence-binding
dimension of this: not whether the intermediate states are accurate, but whether each state
is traceable to the observations that produced it. TRACE~\cite{arxiv240611460} demonstrates
that constructing knowledge-grounded reasoning chains from structured context eliminates
irrelevant retrieval noise in multi-hop queries; the orcast step-log is this structure applied
to skill-dispatch interactions. Chain-of-Thought evaluation~\cite{arxiv220111903} confirms
that intermediate reasoning steps are causally important for correct final outputs, supporting
the claim that step-log quality measurement is foundational to planning system evaluation.

\subsection{Applying $\Runcited$ to planning systems}

The benchmark methodology in Section~4 is directly portable to any AI planning system whose
interactions produce a structured trace. For a planning system that receives a query about a
physical or scientific domain, dispatches skills or sensors to gather current-state information,
produces a plan or action sequence based on gathered context, and narrates the plan with
references to gathered artifacts: inject the planning trace as RAG context to a language model,
ask a domain-specific question, and measure $\Runcited$ on the response. When the trace binds
the query to artifact references, $\Runcited$ approaches 0\%. When the trace is sparse or
unstructured, $\Runcited$ remains high.

The RAG lift hierarchy (Section~5) provides the diagnostic: step-log context achieves full
grounding; structured domain data achieves partial lift; narrow structured data achieves no
lift. This hierarchy is a property of context-type alignment with the science claims the
model would otherwise generate, not a property of context volume. The Ragas framework
~\cite{arxiv230915217} measures faithfulness separately from answer relevancy; $\Runcited$
measures a simpler, deployable condition that precedes both: whether claims are cited at all.
